\section{Problem Setting}
\label{sec:problem_setting}

We address a supervised binary classification problem. Let $\mathbf{x} \in \mathcal{X}$ denote a structured feature vector describing one individual, and let $y \in \{0,1\}$ denote the income label, where $1$ corresponds to annual income greater than \$50K and $0$ otherwise. The goal is to learn a classifier $f(\mathbf{x})$ that predicts the income class for unseen examples.

The Adult dataset contains 48,842 rows with 14 predictive features and one binary income label~\cite{uciadult}. Among the 14 features, 6 are numeric and 8 are categorical. The class distribution is moderately imbalanced: 11,687 positive examples ($>$50K) and 37,155 negative examples ($\leq$50K). The features include age, work class, education, occupation, relationship, sex, capital gain, capital loss, and hours worked per week. This makes the dataset a realistic tabular classification benchmark with mixed data types and non-trivial decision boundaries.

The preprocessing strategy follows the assessment brief. Missing values represented by ``?'' are treated as unavailable observations. Numeric features are imputed with the median, categorical features are imputed with the most frequent category, and categorical variables are transformed through one-hot encoding. Standardization is applied only for the MLP because neural networks are more sensitive to feature scale, whereas tree-based models are generally insensitive to monotonic feature scaling. No target information from the validation or test split is used during preprocessing, which helps avoid data leakage.

This setting is appropriate for the assignment because it uses a public dataset, a clearly defined objective, and a task where both performance and computational trade-offs can be analysed rigorously.
